\scanpage{607}\mspage{16}
% archive sheet number: 606

et \deleted{\unsure{on}} il est clair (comme les $J\in\Phi$ sont
$\ne\emptyset$) que \deleted{pour} les $k_J$ en question sont dans
$\underline S$. Ainsi, pour un filtre $\Phi$ comme dans le cor. 1,
on trouve un système inductif d'objets de $\underline S$, indexé
par $\Phi^{\unsure{o|\circ}}$, dont la limite inductive n'est pas dans
$\underline S$. Donc le a) de la proposition 4 se déduit en notant
que pour tt cardinal $\pi$\inserted{$\in\mathcal U$}, on peut trouver
un ens $E\in\mathcal U$, et un filtre \inserted{$\Phi$} sur $E$ plus
fin que $\Phi_0$, tel que l'ens ordonné $\Phi^o$ soit \unsure{cord|croit|grand}
devant $\pi$. Pour \deleted{\unsure{t}}trouver un tel $\Phi$, il suffit de prendre $E$ de
cardinal $\pi'>\pi$, et pour $\Phi$ l'ens des complémentaires des
parties de $E$ de cardinal $\leq\pi$, de sorte que
$\unsure{\Phi^o|\Phi_0}$ est isom. à l'ens ordonné des
\deleted{\unsure{classes}} parties \inserted{de $E$} de cardinal
$\leq\unsure{\pi|\chi}$, lequel est bien \unsure{gros} devant $\pi$.
Cela prouve \deleted{\unsure{facilt}} le a) de prop. 4.

Prouvons maintenant le b). \unsure{Admettons} le

\uline{Lemme} Soit $M$ une catégorie accessible stable par
$\varinjlim$ \deleted{filtrantes dénombrables}
\inserted{de type $I_\alpha$, $\alpha$ un ordinal limite donné},
$\underline S'$ une sous-catégorie \uline{fermée}, $\pi$ un cardinal utile
